% Pursuing Stacks (cotes 134-2 to 134-8): archive cote 134-5, site page 102 (PDF page 103), typescript page 289.
% typescript. Draft compared: ps4.tex:5189-5310.
% AI-transcribed by Claude Opus 5.5 (claude-opus-5-5, adaptive thinking, effort high), Message Batches API, batch ps2 (msgbatch_01NVwr1Yuo8qJf5rykeGBXPK), 2026-11-12; not yet checked by eye.
\scanpage{134-5/102}
\mspage{289}

\uline{Comments}. In terms of what is known to us since sections 81 and 82,
and notably that any two sections of a 0-connected object of $M$ \unsure{or} of $M'$
are homotopic, the Scholie is just a long tautology. I have taken
great care\inserted{however} in stating it , so as to get as clear a view as possible of
exactly what the relevant relationships are. In a) the conditions
$(*)$ and $(*')$ are just different ways of stating that $f^*$ is a morphism
of contractibility structures, which can still be expressed in various
other ways, compare p.251-252. Similarly, if we assume that $f_!$ com-%
mutes with finite products, then (!) or (!') can be viewed as expres-%
sing that $f_!$ is a morphism of homotopy structures (in opposite directi-%
on), which again could be expressed in various other ways, for instance
in terms of the homotopy relation for maps, or in terms of homotopisms.
This implies that in any case $f^*$ induces a \xout{map for the} functor between
the strict localizations
\[ \overline{f^*}^{c}_{\deleted{\illegible{1}}} : W_c^{-1}M = \overline{M} \hand{\longrightarrow} W'^{-1}_cM' = \overline{M}' \quad , \]
and similarly for $f_!$ when\inserted{\deleted{\illegible{1}}} $f_!$ commutes to finite products. I doubt
however that the latter\inserted{localized} map is of geometric relevance, except\inserted{maybe} in the
cases when \xout{\illegible{1}} $f_!$ gives rise to relations similar to $(\underline{W})$ and $(\underline{W}')$ in
b) above, and both $f_!$ and $f^*$ are model-preserving with respect to
$\underline{W}$-equivalences, and define quasi-inverse equivalences between the
localizations $\mathrm{Hot}_{M,\underline{W}}$ and $\mathrm{Hot}_{M',\underline{W}}$ - a rather exceptional case indeed.

After stating the Scholie, there is scarcely a doubt left in my
mind about the notion of bimorphism which finally peeled out of our
reflections being a relevant one. This \over{s}{S}cholie\hand{,} \deleted{now,} rather than
the ``seducing'' proposition of section 84 (p.275)\hand{,} now seems to me the
adequate \hand{``}answer\hand{''} of the not-so-silly-after-all question of section 46\deleted{.}\inserted{(cf page 95 -- nearly 200 pages ago!).}
The only minor uncertainty remaining in my mind is whether or not in
the notion of a bimorphism for contractibility structures, we should
insist that $f_!$ should commute to finite products. It seems that it
will be hard to check condition (!), except via an apparently weaker
form such as (!'), and using commutation of $f_!$ to products. But on the
other hand the Scholie makes a good sense, without assuming such commu-%
tation property. If the notion of \xout{a} bimorphism is going to be useful,
time\inserted{only} will tell which terminology use is the best.

More puzzling \xout{is} however is this facit, that we still have not
been able to give a satisfactory definition of a \uline{morphism} of \uline{asphericity}
structures, independently of any restrictive assumption on these, such

\margin{102}
